% ============================================================
%  ch02d：§2.4 Frobenius 与 §2.5 Xue 定理（原文 pp.12–16）
% ============================================================
\subsection{Frobenius}

\pnum{}我们现在讨论纳入 Frobenius 之后会发生什么。

现在假设基域 $k$ 是 $\Fb$。假设 $X$ 定义在 $\F_q$ 上；由于 $G$ 是先验定义在 $\ZZ$ 上的，可知 $\Bun_G$ 也自然地定义在 $\F_q$ 上。回忆有理结构可以由 $X$（相应地，$\Bun_G$）的几何 Frobenius 自同态来编码。

因此，存在 Frobenius 自同态（即沿几何 Frobenius 的拉回）作用在 $\Shv(X)$、$\qLisse(X)$ 与 $\Shv(\Bun_G)$ 上。

按 $\LS^{\restr}_{H}$ 的定义，$\qLisse(X)$ 的 Frobenius 自同态诱导 $\LS^{\restr}_{H}$ 的一个“Frobenius”自同态。

\begin{example}
假设 $H = \mathbb{G}_a$。由 §2.2.6，$\LS^{\restr}_{\mathbb{G}_a}$ 是链复形 $C_{\et}(X)[1]$ 的几何化身；这个复形带有其标准的 Frobenius 自同态，而在这本字典下两者自动匹配。

一般地，$\LS^{\restr}_{H}$ 上的 Frobenius 因此可以被看作 $C_{\et}(X)$ 上 Frobenius 的非线性类比。
\end{example}

\pnum{}我们把 $\LS^{\arthm}_{H}$ 定义为 $\LS^{\restr}_{H}$ 的 Frobenius 不动点。更确切地说，我们有形式导出叠的笛卡尔图：
\begin{equation*}
\begin{tikzcd}[row sep=large, column sep=large]
  \LS^{\arthm}_H \arrow[r] \arrow[d] & \LS^{\restr}_H \arrow[d] \\
  \LS^{\restr}_H \arrow[r, "\Delta"] & \LS^{\restr}_H \times \LS^{\restr}_H
\end{tikzcd}
\end{equation*}
其中右端的箭头是 Frobenius 映射的图。

注意 $\LS^{\arthm}_{GL_n}$ 的 $\e$-点是 $X$ 上秩 $n$ 的 lisse Weil 层。更一般地，我们把 $\LS^{\arthm}_{\bG}$ 看作参数化从 Weil 群 $\mathcal{W}_X$ 到代数群 $\bG/\e$ 的连续同态的叠，并把这些同态模共轭。更简洁地说：$\LS^{\arthm}_{\bG}$ 是整体域 $\F_q(X)$ 的无分歧 Langlands 参数的叠。

\begin{remark}
令 $\qLisse^{\phi}(X)$ 表示\emph{拟 lisse Weil 层}的 DG 范畴，按定义它是 $\qLisse(X)$ 上来自 Frobenius 的 $\ZZ$-作用的不动点。同义反复地，$\LS^{\arthm}_{H}$ 参数化对称幺半函子 $\Rep(H) \to \qLisse^{\phi}(X)$，其方式与 $\LS^{\restr}_{H}$ 参数化对称幺半函子 $\Rep(H) \to \qLisse(X)$ 相同（即 $\LS^{\arthm}_{H}$ 的 $S$-点是右 $t$-正合对称幺半函子 $\Rep(H) \to \qLisse(X) \otimes \QCoh(S)$）。

然而，$\qLisse^{\phi}(X)$ 与 $\qLisse(X)$ 具有不同的范畴性质。例如，$\qLisse^{\phi}(X)^{\heartsuit}$ 不是 Tannaka 范畴。这导致两种设定之间的一些形式差异：$\LS^{\arthm}_{H}$ 在某些方面表现得更像 Betti 局部系统的模空间；例如，结果表明（[2] 定理 16.1.4）$\LS^{\arthm}_{H}$ 是一个（非形式的！）代数叠，并且是拟紧的（特别地：只有有限多个连通分支！）。
\end{remark}

\begin{remark}
除了一句评论之外，我们在本节不试图给出更显式的图像：$\LS^{\arthm}_{\bG}$ 的几何比其受限对应物更复杂。但在定理 4.3.3.1 中，我们将在包含平凡表示的一块 $\LS^{\arthm}_{\bG}$ 上给出坐标，从而对其几何提供一些显式分析。
\end{remark}

\pnum{}本质上由注 2.3.1.2，任何受限几何朗兰兹等价\emph{必须}与两边的 Frobenius 自同态相容。

回忆 §2.1.6 中提到的\emph{范畴迹}概念：它取（可对偶化的）DG 范畴连同自函子，产生向量空间。

于是我们可以对受限几何朗兰兹等价的两边取 Frobenius 的迹。如 [2] §16 所概述，$\IndCoh_{\Nilp^{\mathrm{spec}}}(\LS^{\restr}_{\bG})$ 上 Frobenius 的迹与 $\IndCoh(\LS^{\restr}_{\bG})$ 上的相同，即：
\[
  \Gamma(\LS^{\arthm}_{\bG}, \omega),
\]
其中 $\omega$ 是 $\LS^{\arthm}_{\bG}$ 上的\emph{对偶化层}。

另一方面，[4] 的主定理把 $\Shv_{\Nilp}(\Bun_G)$ 上 Frobenius 的迹计算为：
\[
  \Aut^{\unr}_{G,c}.
\]
利用 (2.2)，人们可以把这解释为层—函数对应的一个更高范畴版本（尽管是在特殊情形中，而不是一般的几何现象）。

\pnum{}我们最终得到\emph{算术}猜想：
\[
  \Aut^{\unr}_{G,c} \simeq \Gamma(\LS^{\arthm}_{\bG}, \omega).
\]
正如引言中那样，左端的向量空间是无分歧自守函数的空间，即 $\Bun_G(\F_q)$ 上的紧支撑函数。特别地，无分歧尖形式位于这个空间中。

\pnum{}存在一个典范映射 $\tau: \mathcal{O}_{\LS^{\arthm}_{\bG}} \to \omega_{\LS^{\arthm}_{\bG}}$，编码 $\LS^{\arthm}_{\bG}$ 上的一个“弱 Calabi--Yau”结构——更多讨论参见 §4.6.5。

特别地，存在一个从 $\LS^{\arthm}_{\bG}$ 上的函数到上面右端 $\Gamma(\LS^{\arthm}_{\bG}, \omega)$ 的自然映射。人们应当把 Arthur 的 $SL_2$ 看作在 $\LS^{\arthm}_{\bG}$ 上度量 $\mathcal{O}$ 与 $\omega$ 之间的差异。

\begin{example}
假设 $\sigma \in \LS^{\arthm}_{\bG}$ 是这个叠的一个光滑孤立点。（这样的 $\sigma$ 称为\emph{椭圆}或\emph{离散} Langlands 参数。）于是人们可以看出 $\tau|_\sigma$ 是一个同构。因此，我们的猜想预言存在一个与 $\sigma$ 对应的\footnote{我们对“对应于”的含义有些马虎。更确切地说，我们的猜想结合 [2] §24.2 的讨论，蕴含应当存在 V.~Lafforgue 的 \emph{excursion 代数}作用的一个 $1$ 维无分歧本征形式空间，其特征值由 $\sigma$ 定义。众所周知，对一般 $G$，仅靠 Hecke 算子不足以挑出一个 $1$ 维本征空间。} $1$ 维无分歧自同构形式空间。
\end{example}

\subsection{Xue 的定理与 Frobenius 的迹}

\pnum{}上面我们说：
\begin{equation}
  \mathrm{tr}(\Frob^*_{\Bun_G}, \Shv_{\Nilp}(\Bun_G)) \simeq \Aut^{\unr}_{G,c}
\end{equation}
是 [4] 的主定理。我们简要说明这是如何证明的。关键作用由 Xue 来自 [31] 的定理扮演。

人们也可以转向 [4] 的引言以了解论证概述。我们的概述在这里并无太大不同，只是我们稍微更努力地把 Beilinson 的谱投影扫到地毯下（也许对讨论有害）。

\pnumt{第 1 步。}{} Drinfeld 关于朗兰兹对应的工作\footnote{这里讨论的关于一般约化群的观点来自 [29]。}中的一个教训是：把自守函数 $\Aut^{\unr}_G$ 看作 \emph{shtuka 上同调层}的特殊情形通常很有帮助。

我们简要回顾这个故事。shtuka 构造取一个有限集 $I$、一个表示 $V \in \Rep(\bG^I)$，并产生 $X^I$ 上的一个层 $\Sht_{I,V} \in \Shv(X^I)$。确切地说，附着于 $I$ 与 $V$ 的数据，有一个 \emph{Hecke 函子}：
\[
  H_V: \Shv(\Bun_G) \to \Shv(\Bun_G \times X^I).
\]
这个函子来自一个自然定义的\emph{核} $\mathcal{K}_V \in \Shv(\Bun_G \times \Bun_G \times X^I)$。我们提醒，几何 Satake 在构造中起关键作用，关于构造的更多细节我们参考 [22] 或 [4]。

于是 $\Sht_{I,V}$ 通过沿 Frobenius 的图作 $*$-拉回得到：
\[
  \Bun_G \times X^I \xrightarrow{\mathrm{Graph}_{\Frob} \times \mathrm{id}_{X^I}} \Bun_G \times \Bun_G \times X^I
\]
然后沿 $\Bun_G$ 因子取紧支撑上同调，即推前到 $X^I$。

例如，当 $I = \varnothing$（且 $V$ 是平凡表示的 $1$ 维表示）时，$\mathcal{K}_V = \Delta_!(\mathbf{e}_{\Bun_G})$，所以由基变换，上面的计算给出 $C_{\et,c}(\Bun_G(\F_q)) = \Aut^{\unr}_G$。

shtuka 上同调层之间存在自然的态射。首先，对固定的 $I$，上面的构造给出一个函子 $\Sht_I: \Rep(\bG^I) \to \Shv(X^I)$。但我们也可以变化 $I$；更确切地说，$\Rep(\bG)$ 上的对称幺半结构把指派 $I \mapsto \Rep(\bG^I) = \Rep(\bG)^{\otimes I}$ 变成一个函子 $\mathbf{fSet} \to \DGCat_{\mathrm{cont}}$（其中 $\mathbf{fSet}$ 是有限集的范畴）；沿对角态射的 $*$-拉回使指派 $I \mapsto \Shv(X^I)$ 也变成一个函子 $\mathbf{fSet} \to \DGCat_{\mathrm{cont}}$。于是几何 Satake 的标准函子性说，我们有一个函子之间的自然变换：
\[
  \Sht: (I \mapsto \Rep(\bG^I)) \to (I \mapsto \Shv(X^I)),
\]
即
\[
  \mathbf{fSet} \to \DGCat_{\mathrm{cont}}.
\]
这种函子性是 shtuka 上同调层的一个关键性质，它的存在编码了自守函数的关键对称性：V.~Lafforgue 在 [22] 中正是利用这种函子性来构造 excursion 算子。

\pnumt{第 2 步。}{} 我们现在类似地推广定理的另一边，我们提醒它是 $\mathrm{tr}_{\Shv_{\Nilp}(\Bun_G)}(\Frob)$。答案应当输入 $V \in \Rep(\bG^I)$，并在 $X^I$ 上产生一个层，我们最终将之记为 $\Sht^{\mathrm{tr}}_{I,V}$。当然，构造应当涉及 Hecke 函子与 $\Shv_{\Nilp}(\Bun_G)$，所以我们暂且离题讨论后者。

遵循拓扑设定中的 [27]，我们在 [2] 定理 14.2.4（及其后续讨论）中证明：对 $V \in \Rep(\bG^I)$，Hecke 函子 $H_V$ 把 $\Shv_{\Nilp}(\Bun_G) \subseteq \Shv(\Bun_G)$ 映入 $\Shv_{\Nilp}(\Bun_G) \otimes \qLisse(X^I) \subseteq \Shv(\Bun_G \times X^I)$。

此外，我们在\emph{同上}定理 14.4.3 中证明了一个逆命题：对 $\mathcal{F} \in \Shv(\Bun_G)$，若对所有 $V \in \Rep(\bG)$ 都有 $H_V(\mathcal{F}) \in \Shv(\Bun_G) \otimes \qLisse(X)$，则必有 $\mathcal{F} \in \Shv_{\Nilp}(\Bun_G)$。\footnote{技术上，这个断言中对基域的特征有一些次要限制。回忆 §1.4，我们总是忽略这些小特征，并在讨论中隐含假定我们远离这些特征。}

于是我们用一句格言来总结：$\Shv_{\Nilp}(\Bun_G)$ 可以看作层 $\mathcal{F}$ 的子范畴，使得当我们变化点 $x \in X$ 时，其 \emph{Hecke 变换} $H_{V,x}(\mathcal{F})$ 是局部常值的。

这种关于 $\Shv_{\Nilp}(\Bun_G)$ 的观点实际上对 AGKRRV 系列中几乎\footnote{主要的例外是 [2] 中的 \emph{Künneth 公式}。}每个结果都更好。（从一种观点看，这就是为什么在我们当前的分析中引入 Hecke 函子与一般的 shtuka 层很重要：$\Shv_{\Nilp}(\Bun_G)$ 本身最好用 Hecke 作用来理解。）

\pnumt{第 3 步。}{} 由上，对 $V \in \Rep(\bG)$，我们有一个 Hecke 函子：
\[
  H_V: \Shv_{\Nilp}(\Bun_G) \to \Shv_{\Nilp}(\Bun_G) \otimes \qLisse(X^I).
\]
我们可以把这个函子与 $\Bun_G$ 上的 Frobenius 预复合，得到：
\[
  H_V \circ \Frob_{\Bun_G}: \Shv_{\Nilp}(\Bun_G) \to \Shv_{\Nilp}(\Bun_G) \otimes \qLisse(X^I).
\]
然后我们可以沿\footnote{这类似于说：如果有一个线性变换 $W_1 \to W_1 \otimes W_2$（$W_1$ 有限维），我们有一个相应的向量属于 $W_1^\vee \otimes W_1 \otimes W_2$，并且可以沿前两个因子配对，得到一个“沿 $W_1$ 的迹”，它是 $W_2$ 中的向量。} $\Shv_{\Nilp}(\Bun_G)$ 取迹，得到 $\qLisse(X^I)$ 中的一个对象。这就是所需的对象 $\Sht^{\mathrm{tr}}_{I,V}$。

我们在下面想要证明的是，我们有函子性的等同：
\begin{equation}
  \Sht^{\mathrm{tr}}_{I,V} \simeq \Sht_{I,V}[2|I|],
\end{equation}
其中上同调平移的出现有技术原因，将在下面显现。$I = \varnothing$、$V$ 为 $1$ 维的情形用简洁记号重新得到了 (2.5)。

\pnumt{第 4 步。}{} 观察 $\Sht^{\mathrm{tr}}_{I,V}$ 与 $\Sht_{I,V}$ 之间的一个差异：出于本质的几何原因，$\Sht^{\mathrm{tr}}$ 取值在 $\qLisse(X^I) \subseteq \Shv(X^I)$ 中，而这对 $\Sht_I$ 本身并不显然。

在 [3] 中，我们引入了计算 $\Shv_{\Nilp}(\Bun_G)$ 上迹的方法。细节我们参考\emph{同上}，但简短的答案是：迹可以用\emph{一般的几何要素}（upper-$*$ 与 lower-$!$ 函子）加上\emph{来自（几何）表示论的一个具体要素}来计算。后者是 Beilinson 的谱投影，其工作（就我们的目的而言）是取一个相容的（对 $I$）函子系统 $\mathcal{S}_I: \Rep(\bG^I) \to \Shv(X^I)$，并产生一个相容的系统 $\lambda \mathcal{S}_I: \Rep(\bG^I) \to \qLisse(X^I)$。

在不精确的、高层次的看法下，[3] 的配方给出如下答案：函子 $\Sht^{\mathrm{tr}}$ 的系统是函子 $\Sht_I$ 的系统在取值于 $\qLisse(X^I)$（而不是 $\Shv(X^I)$）中的最佳逼近，即它是 $\lambda \Sht_I$。

于是 Xue 的定理 [31] 说 $\Sht_I$ 本身取值在 $\qLisse(X^I)$ 中，所以 $\lambda \Sht_I$ 与 $\Sht_I$ 本身重合，从而我们得到 (2.6)。

\pnumt{第 5 步。}{} 上面在道德上正确，但我们现在要修正一个谎言。下面的讨论可以与 [4] 注 3.2.6 相比较。

当有限集 $I$ 变化时，函子 $\Sht_I$ 在 upper-$*$ 函子下相容。然而，应用谱投影的程序适用于在 upper-$!$ 函子下相容的函子系统 $\mathcal{S}_I$。所以我们需要 shtuka 函子的一个变体 $\Sht^!$，它在 upper-$!$ 函子下适当相容。

相关的函子 $\Sht^!_{I,-}: \Rep(\bG^I) \to \Shv(X^I)$ 由下面的公式刻画：
\begin{equation}
  C_c(X^I, \Sht^!_{I,V} \boxtimes \mathcal{F}) = C_c(\Bun_G \times X^I, (\Frob \times \mathrm{id})^*(\mathcal{K}_V \boxtimes p_3^*(\mathcal{F})))
\end{equation}
其中 $p_3: \Bun_G \times \Bun_G \times X^I \to X^I$ 是投影。\footnote{一个记号上的说明：我们的记号与 [4] 不一致。我们现在称为 $\Sht^!$ 的函子族被整齐地打包在 [4] 中称为 $\Sht$ 的单个函子里，尽管我们这里称为 $\Sht_I$ 的函子在 [4] 中记号相同。由 Xue 的定理，它们只相差平移。}

先验地，[3] 的结果\emph{实际上}蕴含 $\{\Sht^{\mathrm{tr}}_I\}_{I \in \mathbf{fSet}}$ 是函子 $\{\Sht^!_I\}_{I \in \mathbf{fSet}}$ 在子范畴 $\qLisse \subseteq \Shv$ 中取值的最佳逼近。

逻辑接着通过两次应用 Xue 的定理来进行。首先，该定理告诉我们 $\Sht^!_{I,V} = \Sht_{I,V}[2|I|]$（其中 $2|I|$ 表现为 $2\dim(X^I)$）——我们只需在 (2.7) 的左端用 $\Sht_{I,V}[2|I|]$ 替换 $\Sht^!_{I,V}$，然后利用恒等式 $\mathcal{G} \boxtimes \mathcal{F} = \mathcal{G} \boxtimes \mathcal{F}[-2\dim]$（对 $\mathcal{G}$ lisse）把所得表达式（函子性地）整理成 (2.7) 的右端。特别地，Xue 的定理蕴含 $\Sht^!_{I,V}$ 取值在 $\qLisse$ 中，所以前一段蕴含 $\Sht^{\mathrm{tr}} = \Sht^!$，而我们刚看到这也等于 $\Sht_I[2|I|]$，正如所愿。
